%%
%% This is file `t-markdown.tex',
%% generated with the docstrip utility.
%%
%% The original source files were:
%%
%% markdown.dtx  (with options: `context')
%% 
%% Copyright (C) 2016-2024 Vít Starý Novotný
%% 
%% This work may be distributed and/or modified under the
%% conditions of the LaTeX Project Public License, either version 1.3c
%% of this license or (at your option) any later version.
%% The latest version of this license is in
%% 
%%    http://www.latex-project.org/lppl.txt
%% 
%% This work has the LPPL maintenance status `maintained'.
%% The Current Maintainer of this work is Vít Starý Novotný.
%% 
%% Send bug reports, requests for additions and questions
%% either to the GitHub issue tracker at
%% 
%%   https://github.com/Witiko/markdown/issues
%% 
%% or to the e-mail address <witiko@mail.muni.cz>.
%% 
%% MODIFICATION ADVICE:
%% 
%% If you want to customize this file, it is best to make a copy of
%% the source file(s) from which it was produced.  Use a different
%% name for your copy(ies) and modify the copy(ies); this will ensure
%% that your modifications do not get overwritten when you install a
%% new release of the standard system.  You should also ensure that
%% your modified source file does not generate any modified file with
%% the same name as a standard file.
%% 
%% You will also need to produce your own, suitably named, .ins file to
%% control the generation of files from your source file; this file
%% should contain your own preambles for the files it generates, not
%% those in the standard .ins files.
%% 
%% The names of the source files used are shown above.
%% 
\unprotect
\ifx\ExplSyntaxOn\undefined
  \input expl3-generic
\fi
\ExplSyntaxOn
\tl_const:Nn \c__markdown_option_layer_context_tl { context }
\cs_generate_variant:Nn
  \tl_const:Nn
  { NV }
\tl_if_exist:NF
  \c__markdown_top_layer_tl
  {
    \tl_const:NV
      \c__markdown_top_layer_tl
      \c__markdown_option_layer_context_tl
  }
\ExplSyntaxOff
\writestatus{loading}{ConTeXt User Module / markdown}%
\startmodule[markdown]
\def\dospecials{\do\ \do\\\do\{\do\}\do\$\do\&%
  \do\#\do\^\do\_\do\%\do\~}%
\input markdown/markdown
\let\startmarkdown\relax
\let\stopmarkdown\relax
\let\inputmarkdown\relax
\ExplSyntaxOn
\cs_new:Npn
  \setupmarkdown
  [ #1 ]
  {
    \__markdown_setup:n
      { #1 }
  }
\cs_new:Nn \__markdown_caseless:N
  {
    \regex_replace_all:nnN
      { ([a-z])([A-Z]) }
      { \1 \c { str_lowercase:n } \cB\{ \2 \cE\} }
      #1
    \tl_set:Nx
      #1
      { #1 }
  }
\seq_gput_right:Nn \g__markdown_cases_seq { __markdown_caseless:N }
\str_if_eq:VVT
  \c__markdown_top_layer_tl
  \c__markdown_option_layer_context_tl
  {
    \__markdown_define_option_commands_and_keyvals:
    \__markdown_define_renderers:
    \__markdown_define_renderer_prototypes:
  }
\ExplSyntaxOff
\def\markdownMakeOther{%
  \count0=128\relax
  \loop
    \catcode\count0=11\relax
    \advance\count0 by 1\relax
  \ifnum\count0<256\repeat
  \catcode`|=12}%
\long\def\inputmarkdown{%
  \dosingleempty
  \doinputmarkdown}%
\long\def\doinputmarkdown[#1]#2{%
  \begingroup
    \iffirstargument
      \setupmarkdown[#1]%
    \fi
    \markdownInput{#2}%
  \endgroup}%
\startluacode
  document.markdown_buffering = false
  local function preserve_trailing_spaces(line)
    if document.markdown_buffering then
      line = line:gsub("[ \t][ \t]$", "\t\t")
    end
    return line
  end
  resolvers.installinputlinehandler(preserve_trailing_spaces)
\stopluacode
\begingroup
  \catcode`\|=0%
  \catcode`\\=12%
  |gdef|startmarkdown{%
    |ctxlua{document.markdown_buffering = true}%
    |markdownReadAndConvert{\stopmarkdown}%
                           {|stopmarkdown}}%
  |gdef|stopmarkdown{%
    |ctxlua{document.markdown_buffering = false}%
    |markdownEnd}%
|endgroup
\ExplSyntaxOn
\cs_gset:Nn
  \__markdown_load_theme:nn
  {
    \file_if_exist:nTF
      { t - markdown theme #2.tex }
      {
        \msg_info:nnn
          { markdown }
          { loading-context-theme }
          { #1 }
        \usemodule
          [ t ]
          [ markdown theme #2 ]
      }
      {
        \__markdown_plain_tex_load_theme:nn
          { #1 }
          { #2 }
      }
  }
\msg_new:nnn
  { markdown }
  { loading-context-theme }
  { Loading~ConTeXt~Markdown~theme~#1 }
\ExplSyntaxOff
\markdownIfOption{noDefaults}{}{
  \setupmarkdown[theme=witiko/markdown/defaults]
}
\stopmodule
\protect
\endinput
%%
%% End of file `t-markdown.tex'.
